% ============================================================
%  第三讲
% ============================================================
\section{第三讲}

\subsection{系数（续）}
\srcpage{8}

上次我们定义了 \(\coeff_0\colon\Shv(\Bun_G)\to\Vect\)。现在对 \(X\) 上任意有效除子 \(D\) 定义 \(\coeff_D\)。

考虑叠
\[
\Bun^{\Omega(-D)}_N:=\Bun_B\times_{\Bun_{\Gm}}\{\Omega(-D)\}.
\]
它分类扩张
\[
[0\to\Omega(-D)\to\cE\to\cO\to 0].
\]
设
\[
p_{N,D}\colon\Bun^{\Omega(-D)}_N\to\Bun_G
\]
是忘掉扩张结构、只记住 \(\cE\) 的映射。考虑复合
\[
\phi_D\colon\Bun^{\Omega(-D)}_N\to H^1(X,\Omega(-D))\to H^1(X,\Omega)\simeq\A^1,
\]
其中最后一个同构来自 Serre 对偶。

\begin{definition}\label{def:coeffD}
定义
\[
\coeff_D(F):=\mathrm{C}^\bullet_c\bigl(\Bun^{\Omega(-D)}_N,\ p_{N,D}^!(F)\overset{!}{\otimes}\phi_D^!(\exp)\bigr).
\]
\end{definition}

现在可以陈述我们的目标了。

\begin{goal}\label{goal:main}
我们要作如下构造：

\begin{itemize}
\item （输入）\(X\) 上一个不可约 \(\mathrm{SL}_2\)-局部系统 \(\sigma\)；
\item （输出）一个尖点反常层 \(\cF_\sigma\in\Shv(\Bun_{\mathrm{PGL}_2})\)，它在
\(\Bun^{\mathrm{even}/\mathrm{odd}}_{\mathrm{PGL}_2}\) 上不可约，并附带以下数据：
\begin{itemize}
\item 对每个有效除子 \(D=\sum n_i x_i\)，一个同构
\begin{equation}\label{eq:3.1}
\coeff_D(\cF_\sigma)\simeq\bigotimes_i\Sym^{n_i}(\sigma_{x_i}).
\end{equation}
\end{itemize}
\end{itemize}
\end{goal}

\begin{remark}\label{rem:recursion}
\eqref{eq:3.1} 类比于上一讲的乘法性与递推公式 \eqref{eq:2.1}，因为对 \(\mathrm{SL}_2\) 的标准表示 \(V\)，有
\[
\Sym^n(V)\otimes V\simeq\Sym^{n+1}(V)\oplus\Sym^{n-1}(V).
\]
\end{remark}

上述同构 \eqref{eq:3.1} 应当代数地依赖于 \(D\)。更精确地说，对每个 \(d\ge 0\)，可以构造函子
\[
\coeff_d\colon\Shv(\Bun_G)\to\Shv(\Sym^d(X)),
\]
使得对 \(D\in\Sym^d(X)\)，\(\coeff_D\) 可以等同于复合
\[
\Shv(\Bun_G)\xrightarrow{\ \coeff_d\ }\Shv(\Sym^d(X))\xrightarrow{\ (-)|^!_D\ }\Vect.
\]
于是 \eqref{eq:3.1} 应当升级为同构
\[
\coeff_d(\cF_\sigma)\simeq\sigma^{(d)},
\]
其中
\[
\sigma^{(d)}:=\Add_{d,*}(\sigma^{\boxtimes d})^{S_d}
\]
按秩 \(1\) 情形中的方式定义。然而，\(\sigma^{(d)}\) 是反常的，但不是 lisse 的。
\yizhu{原文说 \(\sigma^{(d)}\) ``is perverse but not lisse''，这是笔误：\(\sigma^{(d)}\) 是 \(\Sym^d(X)\)
上的局部系统，按定义就是 lisse 的；这里想说的应当是「\(\cF_\sigma\) 是反常的，但不是 lisse 的」，
与秩 \(1\) 情形中 \(\chi_\sigma\) 为 lisse 形成对照。按反常层的标准归一化，\(\Sym^d(X)\)（\(d\) 维）上的
局部系统要移位 \(d\) 次才落在反常 t-结构的心内；原文断言的是 \(\cF_\sigma\)，不涉及这个移位。}{原文笔误·A：与上一段「不再要求 lisse，只要求反常」对照；B：nLab ``perverse sheaf''}

基于模形式的动机，我们可能希望函子
\[
\prod_{d\ge 0}\coeff_d\colon\Shv(\Bun_G)_{\mathrm{cusp}}\to\prod_{d\ge 0}\Shv(\Sym^d(X))
\]
是全忠实的（从而 \(\cF_\sigma\) 由上述预期唯一确定）。然而这并不成立：右端按 \(\Z_{\ge 0}\) 分裂成若干因子，
而左端只按 \(\Z/2\) 分裂。尽管如此，我们有下述断言。

\begin{proposition}\label{prop:determined}
设 \(F\) 是 \(\Bun_{\mathrm{PGL}_2}\) 上的反常尖点层，且它在
\(\Bun^{\mathrm{even}/\mathrm{odd}}_{\mathrm{PGL}_2}\) 上的限制不可约、支撑稠密。
则对任意 \(d>N\)，\(F\) 由 \(\coeff_d(F)\) 与 \(\coeff_{d+1}(F)\) 唯一确定，其中 \(N\) 是只依赖于 \(g\) 的整数。
\end{proposition}

\begin{remark}[警告]\label{rem:only-pgl2}
上述断言是 \(\mathrm{PGL}_2\) 特有的。
\end{remark}

\subsection{Fourier 变换}
\srcpage{9}

为解释命题 \ref{prop:determined}，我们回忆 Fourier（--Deligne）变换。对有限维向量空间 \(V\) 及其对偶 \(V^\vee\)，
考虑对应
\[
\begin{tikzcd}
& V\times V^\vee \arrow[dl,"\mathrm{pr}_1"'] \arrow[dr,"\mathrm{pr}_2"] & \\
V & & V^\vee
\end{tikzcd}
\]
有等价
\[
\Four\colon\Shv(V)\xrightarrow{\ \simeq\ }\Shv(V^\vee),\qquad
F\mapsto \mathrm{pr}_{2,*}\bigl(\mathrm{pr}_1^!(F)\overset{!}{\otimes}\mathrm{ev}^!(\exp)\bigr),
\]
其中 \(\mathrm{ev}\colon V\times V^\vee\to\A^1\) 是配对映射。这可以推广到基 \(S\) 上的向量丛 \(E\) 与 \(E^\vee\)，
于是有典范等价
\[
\Four\colon\Shv(E)\xrightarrow{\ \simeq\ }\Shv(E^\vee).
\]

\begin{remark}\label{rem:four-betti}
在 \(\exp\) 不存在的 Betti 情形，我们可以对单值（monodromic）层作 Fourier 变换。
\end{remark}

我们可以用 Fourier 变换与丛的几何重写 \(\coeff_d\)。回忆映射
\[
q_{-d}\colon\Bun^{-d}_B\to\Bun^{-d}_{\Gm}
\]
在 \(L\) 处的纤维等同于 \(R\Gamma(X,L)[1]\)。对 \(d>0\)，有 \(H^0(X,L)\simeq 0\)，因此 \(q_{-d}\) 是
\(\Bun^{-d}_{\Gm}\) 上的向量丛，记作 \(E_d\)。对偶丛 \(E_d^\vee\) 在 \(L\) 处的纤维为
\[
H^1(X,L)^\vee\simeq H^0(X,L^\vee\otimes\Omega).
\]
于是 \(E_d^\vee\) 是分类 \((L,s\colon L\to\Omega)\) 的模叠，其中 \(L\in\Bun^{-d}_{\Gm}\)，\(s\colon L\to\Omega\) 是映射。

注意有识别
\[
\Sym^{d+(2g-2)}(X)\to E_d^\vee\setminus\{0\},\qquad D\mapsto[\Omega(-D)\to\Omega],
\]
其中 \(0\) 表示 \(E_d^\vee\) 的零截面（它同构于 \(\Bun^{-d}_{\Gm}\)）。

\begin{exercise}\label{ex:coeff-four}
函子 \(\coeff_{d+(2g-2)}\) 可以等同于复合
\[
\Shv(\Bun_G)\xrightarrow{\ p^!_{-d}\ }\Shv(\Bun^{-d}_B)\xrightarrow{\ \Four\ }\Shv(E_d^\vee)
\xrightarrow{\ \text{限制}\ }\Shv(\Sym^{d+(2g-2)}(X)).
\]
\end{exercise}

\begin{fact}\label{fact:connected}
若 \(f\colon Y\to Z\) 是纤维连通的平滑映射，则反常层上的拉回函子是全忠实的。
\end{fact}

\begin{fact}\label{fact:connected-irr}
若 \(f\colon Y\to Z\) 是纤维连通或空的平滑映射，且 \(Y\) 非空，则拉回函子在支撑稠密的不可约反常层上是全忠实的。
\end{fact}

\begin{exercise}\label{ex:fibers}
对任意开子叠 \(U_n\subset\Bun_{\mathrm{PGL}_2}\)，存在只依赖于 \(n\) 与 \(g\) 的整数 \(N\)，
使得对任意 \(d>N\)，映射 \(p_{-d}\colon\Bun^{-d}_B\to\Bun_G\) 在 \(U_n\) 上的纤维连通或空。
\end{exercise}

\begin{proof}[命题 \ref{prop:determined} 的证明]
回忆 \(F\) 由它在 \(U_n\subset\Bun_{\mathrm{PGL}_2}\) 上的限制决定，其中 \(n\) 是只依赖于 \(g\) 的整数
（推论 \ref{cor:extend}）。于是由上述事实与练习，存在只依赖于 \(g\) 的整数 \(N\)，使得对任意 \(d>N\)，
函子 \(p^!_{-d}\) 不会丢失 \(F|_{\Bun^{d+2\Z/2\Z}_{\mathrm{PGL}_2}}\) 的信息。因此，为证明命题 \ref{prop:determined}，
只需从 \(\Four(p^!_{-d}F)\) 沿开浸入
\[
E_d^\vee\setminus 0\subset E_d^\vee
\]
的限制恢复 \(\Four(p^!_{-d}F)\)。

为此注意：\(\Four(p^!_{-d}F)\) 在零截面上的 \(!\)-限制可以等同于 \(\CT^{-d}_*(F)\)；
由于 \(F\) 被假定为尖点，后者为零。
\yizhu{命题 \ref{prop:determined} 的证明中，「零截面上的 \(!\)-限制为零」是关键：它说明
\(\Four(p^!_{-d}F)\) 在 \(E_d^\vee\) 的零截面附近没有「额外」的部分，因此它由开子集
\(E_d^\vee\setminus 0\) 上的限制唯一决定。这一步对应模形式中「常数项为零」的尖点条件。}{原文省略·I：本工程 §2.1 与 §3.2 的衔接}
\end{proof}

\subsection{几何 Langlands 对应}
\srcpage{10}

我们已经看到：若 \(\cF_\sigma\) 存在，则它由 \(\sigma\) 唯一确定。下述结果称为 \(\mathrm{PGL}_2\) 的几何 Langlands 对应。

\begin{theorem}[Drinfeld, Laumon, Frenkel--Gaitsgory--Vilonen]\label{thm:gl2}
对任意不可约 \(\mathrm{SL}_2\)-局部系统 \(\sigma\)，存在 \(\cF_\sigma\in\Shv(\Bun_{\mathrm{PGL}_2})\)，
它反常，且在 \(\Bun_{\mathrm{PGL}_2}\) 的每个连通分支上不可约，并满足
\[
\coeff_d(\cF_\sigma)\simeq\sigma^{(d)}.
\]
此外，\(\cF_\sigma\) 是 Hecke 特征层。
\end{theorem}

为解释「Hecke 特征层」这个概念，考虑模叠 \(\mathrm{Hecke}\)，它分类 \((x\in X,\cE\subset\cE'\subset\cE(x))\)，
其中（模去张量线丛）\(\cE\) 与 \(\cE'\) 是 \(X\) 上的秩 \(2\) 丛，且包含 \(\cE\subset\cE'\subset\cE(x)\) 是严格的。
有对应
\[
\begin{tikzcd}
& \mathrm{Hecke} \arrow[dl,"\overleftarrow{h}"'] \arrow[dr,"\overrightarrow{h}"] & \\
\Bun_G & & \Bun_G\times X
\end{tikzcd}
\]
其中 \(\overleftarrow{h}\) 把 \((x\in X,\cE\subset\cE'\subset\cE(x))\) 送到 \(\cE\)，
\(\overrightarrow{h}\) 把它送到 \((\cE',x)\)。考虑函子
\yizhu{Hecke 叠的严格包含 \(\cE\subset\cE'\subset\cE(x)\) 可以这样理解：\(\cE'/\cE\) 是 \(x\) 处的
一维「跳跃」，而 \(\cE(x)=\cE\otimes\cO(x)\) 是给 \(\cE\) 在 \(x\) 处加一个极点。模去张量线丛后，
这个三元组只依赖 \((\cE,\cE')\) 的相对位置。Hecke 函子 \(H_{\mathrm{std}}\) 的「先 \(!\) 拉回、再 \(*\) 推前」
是几何 Langlands 中 Hecke 算子的标准形式。}{背景直觉·A：Frenkel--Gaitsgory--Vilonen, \emph{On the geometric Langlands conjecture}, J. Amer. Math. Soc. 15 (2002)；B：nLab ``geometric Langlands correspondence''}
\begin{equation}\label{eq:3.2}
H_{\mathrm{std}}:=\overrightarrow{h}_*\,\overleftarrow{h}^!.
\end{equation}
Hecke 特征性质说的是
\begin{equation}\label{eq:3.3}
H_{\mathrm{std}}(\cF_\sigma)\simeq\cF_\sigma\boxtimes\sigma.
\end{equation}
